% Generated draft: requires mathematical and source review.
\documentclass[11pt]{article}
\usepackage[margin=25mm]{geometry}
\usepackage{fontspec}
\setmainfont{Latin Modern Roman}
\usepackage{amsmath,amssymb,mathtools,unicode-math}
\setmathfont{Latin Modern Math}
\usepackage{enumitem,needspace,xurl,hyperref}
\setlength{\parindent}{0pt}
\setlength{\parskip}{5pt}
\setlength{\emergencystretch}{3em}
\newcommand{\N}{\mathbb{N}}
\newcommand{\Z}{\mathbb{Z}}
\newcommand{\Q}{\mathbb{Q}}
\newcommand{\R}{\mathbb{R}}
\newcommand{\C}{\mathbb{C}}
\newcommand{\F}{\mathbb{F}}
\newcommand{\A}{\mathbb{A}}
\newcommand{\PP}{\mathbb{P}}
\newcommand{\E}{\mathbb{E}}
\newcommand{\eps}{\varepsilon}
\newcommand{\dd}{\,\mathrm{d}}
\begin{document}
\section*{UnsolvedMath ID: 30002702}
\textbf{Source problem code:} catalogue entry
\subsection{Definitions and mathematical statement}
The supplied record asks whether the Kaplan--Yorke equality holds for a prevalent set of coupling functions in bidirectionally coupled skinny baker's maps. It does not specify the maps, their parameter ranges, the coupling rule, the admissible function space, the invariant measure, or the dimension entering the equality. Consequently, the record determines a question schema rather than a fully specified mathematical conjecture.

To express that schema, let \(\mathcal C\) denote the as-yet unspecified space of admissible bidirectional couplings, let \(T_c\) denote the corresponding coupled map for \(c\in\mathcal C\), and let \(\mu_c\) denote the invariant probability measure intended by the source. Write \(D_*(\mu_c)\) for the dimension intended in the source's Kaplan--Yorke equality; the meaning of \(D_*\) is unresolved.

For reference, if a measure has a specified finite list of finite Lyapunov exponents \(\lambda_1\geq\cdots\geq\lambda_d\), counted with multiplicity, put
\[
S_0=0,\qquad S_k=\sum_{i=1}^{k}\lambda_i,\qquad
j=\max\{k\in\{0,\ldots,d\}:S_k\geq0\}.
\]
The Kaplan--Yorke dimension associated with this list is
\[
D_{\mathrm{KY}}=
\begin{cases}
j+\dfrac{S_j}{|\lambda_{j+1}|},&j<d,\\
d,&j=d.
\end{cases}
\]
When \(j<d\), the definition of \(j\) implies \(\lambda_{j+1}<0\). Applying this formula to \(\mu_c\) requires a specified Lyapunov spectrum and hypotheses ensuring that it is well defined; these are not supplied. In particular, no ergodicity or choice among invariant measures is asserted here.

For reference, in a real complete metric topological vector space \(V\), a subset \(P\subseteq V\) is called prevalent if \(V\setminus P\) is contained in a Borel set \(B\) for which there is a compactly supported Borel probability measure \(\nu\) on \(V\) satisfying
\[
\nu(v+B)=0\qquad\text{for every }v\in V.
\]
The source does not identify an ambient space \(V\), its topology, or whether prevalence is intended in the whole space or relative to an admissible parameter domain.

With those choices and all necessary existence hypotheses supplied, the intended question has the form
\[
\text{Is }\mathcal P=\{c\in\mathcal C:
D_*(\mu_c)=D_{\mathrm{KY}}(\mu_c)\}
\text{ prevalent in the intended coupling space?}
\]
This is a conditional formulation only. It does not select an invariant measure, interpret \(D_*\) as information or Hausdorff dimension, impose parameter restrictions, or resolve whether the question is intended for every admissible choice of the underlying skinny baker's maps.
\subsection{Short English statement}
For bidirectionally coupled skinny baker's maps, is the dimension equality predicted by the Lyapunov exponents valid for a prevalent set of coupling functions? The supplied record leaves the maps, parameters, coupling space, invariant measure, and dimension unspecified, so the displayed formulation requires these choices before it becomes a precise conjecture.
\subsection{Sources}
\begin{itemize}
\item Supplied local UnsolvedMath registry record, permanent ID \(30002702\), source problem code \texttt{catalogue entry}, title: ``Kaplan--Yorke Equality for Coupled Baker's Maps.'' The record belongs to the Oberwolfach Reports Open Problems collection and gives proposed year \(2014\); no proposer is named.
\item Supplied report citation: \emph{Mini-Workshop: Differentiable Ergodic Theory, Dimension Theory and Stable Foliations (2015)}, Oberwolfach Reports, pp.~2569--2617, DOI \texttt{10.4171/owr/2014/46}. \url{https://doi.org/10.4171/owr/2014/46}. The supplied background describes Oberwolfach Reports as documenting mathematical workshops and their open-problem discussions.
\item Supplied literature source: Gr\"oger--Hunt, \emph{Coupled skinny baker's maps}. \url{https://arxiv.org/abs/1303.0030}.
\item Supplied literature source: \emph{Oberwolfach Report 46/2014}. \url{https://doi.org/10.4171/owr/2014/46}.
\item Literature assessment dated exactly \texttt{2026-08-21}: the supplied curation labels the problem open and reports that the cited work proves the Kaplan--Yorke equality for one-way coupling in a prevalent set of cases, proves failure in other cases, and leaves the bidirectional-coupling conjecture unresolved. This is supplied dated curation derived from the OWR material, not an independently verified literature assessment. Time-sensitive status and the cited claims require independent verification.
\item The supplied statement status is \texttt{corrected\_verified}; its verification note reports manual semantic, typographic, and LaTeX review of a meaning-preserving cleanup claimed to be self-contained and to add no unsupported source-dependent content. That supplied label does not establish the missing mathematical specifications or constitute independent verification here.
\end{itemize}
\end{document}
